\scanpage{3}
\mspage{I}
\uline{Calcul de} $S_0\mathcal{T}_G$, \uline{pour} $G$ \uline{\unsure{graphe}} MD \uline{de dim. modulaire} $0$.

\margin{\unsure{On notera} plutôt $S_0^{\circ}\mathcal{T}_G$ \unsure{pour ces} groupoïdes \unsure{étriqués|extrêmes}, \unsure{car} \unsure{bien trop petits}\unsure{]|)}}

$S_0\mathcal{T}_G$ est un $\mathbb{Z}^{\widehat{A}}$-groupoïde, déduit du $\mathbb{Z}^{\widehat{\vec A}}$-groupoïde $S_0\mathcal{T}_{G^b}$ par le changement de groupe \deleted{\unsure{en}} \unsure{interne}
\[
\mathbb{Z}^{\widehat{\vec A}}\longrightarrow\mathbb{Z}^{\widehat A}
\qquad\text{associé à}\qquad
\widehat{\vec A}\longrightarrow\widehat A
\]
i.e.
\[
(1)\qquad S_0\mathcal{T}_G=S_0\mathcal{T}_{G^b}\wedge_{\mathbb{Z}^{\widehat{\vec A}}}\mathbb{Z}^{\widehat A}
\]
et
\[
(2)\qquad S_0\mathcal{T}_{G^b}=\prod_{s\in S}S_0\mathcal{T}_{G_s^b}
=\prod_{s\in S}S_0\mathcal{T}_{0,\widehat{\vec A}_s}
\]
ce qui nous \unsure{ramène} à décrire $S_0\mathcal{T}_G$, quand $G$ est un \unsure{démigraphe} ayant un seul sommet, donc can$^{\mathrm t}$. isomorphe à un $G_{0,I}$, $I$ étant un ens. à trois éléments, \deleted{de sorte} -- i.e. on est ramené à décrire \deleted{\illegible{1}} $S_0\mathcal{T}_{0,I}$.

Considérons \unsure{la} \deleted{courbe}\inserted{droite} projective $I$-marquée sur $\mathbb C$, $\mathbb P^1_I$. Soit $\omega(I)$ l'ens des deux orientations de $I$. Si $\alpha\in\omega(I)$, on a des iso. canoniques \inserted{de $\mathbb U$-torseurs} pour les espaces tangents unitaires $T_i$ \unsure{en les} $i\in I\subset\mathbb P^1_I$
\[
T_i\overset{\varphi_{\alpha,i}}{\underset{\sim}{\longleftrightarrow}}\mathbb U
=\{z\in\mathbb C\mid |z|=1\}
\]
i.e. un isom de $\mathbb U^I$-torseurs
\[
\mathbb U^I\xrightarrow[\sim]{\varphi}P\mathbb U_G=P\mathbb U_{0,I}=\prod_{i\in I}T_i
\]
i.e. un point \deleted{noté} $\bar\alpha=\varphi(1)\in P\mathbb U_{0,I}$, d'où un plongement canonique
\[
(3)\qquad\omega(I)\hookrightarrow P\mathbb U_{0,I}
\]
compatible avec l'homomorphisme sur les groupes d'opérateurs

$(4)$\quad
\diagram{\begin{tikzcd}[ampersand replacement=\&, column sep=large,row sep=large]
\{\pm1\}\arrow[rr]\arrow[dr,"\mathrm{diag}"']\&\&\mathbb U^I\\
\&\{\pm1\}^I\arrow[ur,"{\substack{\text{déduit de l'inclusion}\\\{\pm1\}\hookrightarrow\mathbb U}}"']\&
\end{tikzcd}}
\mspage{2}
